% \documentclass[11pt,a4paper,twoside]{article}

\usepackage[utf8]{inputenc}

\usepackage[T1,T2A]{fontenc}

\usepackage[english.ancient,russian.ancient]{babel}

\usepackage{graphicx}

\usepackage{array}

\usepackage[doi]{samgtu-bib}

\usepackage{samgtu}

\usepackage{samgtu-name}

\usepackage{mathtools}

\mathtoolsset{showonlyrefs}

\usepackage{desclist}

\usepackage{euscript}

\usepackage{mathrsfs} 

\usepackage{multicol}

\usepackage{mathrsfs}

\usepackage{cite}

\usepackage{cmap}

\usepackage{qrcode}

\allowdisplaybreaks[4]

\usepackage[nodayofweek]{datetime}

\usepackage[thinlines]{easybmat}



\renewcommand{\JournalYear}{2018}

\renewcommand{\JournalVolume}{22}

\renewcommand{\JournalNumber}{4}



\newdate{JournalDateSign}{29}{12}{2018}% Дата подписания Вестника в печать

\renewcommand{\JournalDateSign}{\displaydate{JournalDateSign}}

\newdate{JournalDatePrint}{16}{01}{2019}% Дата выхода Вестника из печати

\renewcommand{\JournalDatePrint}{\displaydate{JournalDatePrint}}

%\renewcommand{\JournalUPL}{16.56} % Усл. печ. л.

%\renewcommand{\JournalUIL}{16.53} % Уч.-изд. л.

\renewcommand{\JournalUPL}{15.24} % Усл. печ. л.

\renewcommand{\JournalUIL}{15.21} % Уч.-изд. л.

\renewcommand{\JournalRegNumber}{220/18} % Регистрационный номер

\renewcommand{\JournalOrderNumber}{2} % Номер заказа



%\renewcommand{\JournalRMonth}{Март} % Месяц выхода Вестника

%\renewcommand{\JournalEMonth}{March} % Месяц выхода Вестника

%\renewcommand{\JournalRMonth}{Июнь} % Месяц выхода Вестника

%\renewcommand{\JournalEMonth}{June} % Месяц выхода Вестника

%\renewcommand{\JournalRMonth}{Сентябрь} % Месяц выхода Вестника

%\renewcommand{\JournalEMonth}{September} % Месяц выхода Вестника

\renewcommand{\JournalRMonth}{Декабрь} % Месяц выхода Вестника

\renewcommand{\JournalEMonth}{December} % Месяц выхода Вестника



\newcommand{\totop}[1]{\raisebox{6pt}[1pt][2pt]{#1}}

\newcommand{\tottop}[1]{\raisebox{12pt}[1pt][2pt]{#1}} 

\newcommand{\totttop}[1]{\raisebox{18pt}[1pt][2pt]{#1}} 

\newcommand{\tottttop}[1]{\raisebox{24pt}[1pt][2pt]{#1}} 

\newcommand{\totttttop}[1]{\raisebox{30pt}[1pt][2pt]{#1}} 

\newcommand{\tobot}[1]{\raisebox{-6pt}[1pt][2pt]{#1}} 

\newcommand{\tobbot}[1]{\raisebox{-12pt}[1pt][2pt]{#1}} 

\newcommand{\tobbbot}[1]{\raisebox{-18pt}[1pt][2pt]{#1}}



\graphicspath{{./00PIC/}}
%
%
%%\samgtupubl % для генерации pdf, отдаваемого в типографию СамГТУ
%\pdfcrop % обрезка pdf для размещения на math-net.ru
%\draft % На корректуре
%%\filllastpage % Последняя страница не заполнена не полностью
%%\briefcomm % Статья является кратким сообщением
%
\vsgtu{1515}
%
\FirstPage{236}
\LastPage{248}
%%
%%
%\begin{document}




\renewcommand{\baselinestretch}{0.93}

\hfill \qrcode[hyperlink,height=.8in]{http://doi.org/10.14498/vsgtu1515}
\vspace{-.8in}




\udc{517.968.73}%Обработка текста. Подготовка текстов : Языки программирования. Компьютерные языки
\msc{45J05, 34K05, 34K30}% Коды MSC см. http://www.ams.org/msc/

\rutitle{Построение аналога теоремы Фредгольма для~одного~класса модельных интегро"=дифференциальных уравнений первого порядка с~логарифмической особенностью в~ядре}

\rutitleshort{Построение аналога теоремы Фредгольма\dots}
%
%\rutitlehack{}

\entitle{A construction of analog of Fredgolm theorems for one class of first order model integro-differential equation with logarithmic singularity in the kernel}



\entitleshort{A construction of analog of the Fredholm theorem\dots}



\ruauthor{С.~К.~Зарипов}{З\,а\,р\,и\,п\,о\,в~С.~К.}

\ruaffil{\ruaffid{7258}{Таджикский национальный университет,\\
Таджикистан, 734025, Душанбе, пр. Рудаки 17}.}

\enaffil{\enaffid{7258}{Tajik National University, \\
17, av. Rudaky, Dushanbe, 734025, Tajikistan}.}

\ruauthorsdetails{1}{
\fullauthor{\rupid{121694}{Сарвар Кахрамонович Зарипов}} 
\CAuthor \\ % Автор-корреспондент
%\orcid{} % ORCID ID автора 
\regalia{кандидат физико-математических наук} 
\position{старший преподаватель} 
\dept{каф. математического анализа и теории функций} 
\email{sarvar8383@list.ru}
}



\enauthorsdetails{1}{
\fullauthor{\enpid{121694}{Sarvar~K.~Zaripov}} 
\CAuthor\\  % Автор-корреспондент
%\orcid{} % ORCID ID автора 
\regalia{Cand. Phys. \& Math. Sci.} 
\position{Senior Lecturer} 
\dept{Dept. of  Mathematics Analysis and Function Theory} 
\email{sarvar8383@list.ru}
}


\rukeywords{модельное интегро-дифференциальное уравнение, граничные сингулярные точки, интегральные представления, логарифмическая особенность, характеристическое уравнение}

\enkeywords{integro-differential equation, boundary singular points, manifold solution, integral representation, characteristic equation}

\ruabstract{Для одного модельного интегро-дифференциального уравнения первого порядка с логарифмической особенностью в ядре в зависимости от корней характеристического уравнения найдены интегральные представления многообразия решений через произвольные постоянные. Найдены случаи, когда данное интегро-дифференциальное уравнение имеет единственное решение. Построены аналоги теоремы Фредгольма для этого интегро-дифференциального уравнения. Использованный метод можно применять для изучения модельных и немодельных интегро-дифференциальных уравнений высших порядков.}



\enauthor{S.~K.~Zaripov}{Z\,a\,r\,i\,p\,o\,v~S.~R.}


\enabstract{The integral representations of the solution manifold for one class of the first order model integro-differential equation with logarithmic singularity in the kernel are constructed using arbitrary constants. The cases when the given integro-differential equation has unique solution are found. The analogue of Fredholm theorem is built for given integro-differential equation. The method of solving this problem can be used for the solving of higher order model and non-model integro-differential equations with singular coefficients.}

\newdate{zaripova}{20}{10}{2016}
\date{\displaydate{zaripova}}
\newdate{zaripovb}{14}{03}{2017}
\revisiondate{\displaydate{zaripovb}}
\newdate{zaripovc}{12}{06}{2017}
\accepteddate{\displaydate{zaripovc}}

\newdate{zaripove}{09}{07}{2017}
\renewcommand{\JournalDataOnline}{\displaydate{zaripove}}


\makerutitle




 \Section[n]{Введение}
 Многие задачи прикладного характера, например задача Вольтерра о крутильных колебаниях \cite{zar1:bib:Voltera}, задача Прандтля расчёта крыла самолета [\citen{zar2:bib:Magnaradze,zar3:bib:Vekua}], задача об изучении кинетического уравнения Больцмана, приводят к изучению интегро"=дифференциальных уравнений. 
 Существует большое количество работ для изучения обыкновенных интегро"=дифференциальных уравнений с~регулярными коэффициентами [\citen{zar4:bib:Veynberg,zar5:bib:Nekrasov}]. Также в последние годы в силу своей прикладной важности изучаются прямые и обратные задачи для интегро-дифференциальных уравнений, а также разрабатываются методы для приближённого решения этих уравнений [\citen{zar6:bib:Durdiev,zar7:bib:Durdiev,zar8:bib:Safarov,zar9:bib:Yuldoshev,zar10:bib:YuldoshevT,zar11:bib:YuldoshevT}]. 
 Одним из важных разделов теории интегро"=дифференциальных уравнений являются уравнения с~сингулярными коэффициентами [\citen{zar2:bib:Magnaradze,zar3:bib:Vekua,zar12:bib:MagnaradzeG}] или с~сингулярно возмущенными коэффициентами [\citen{zar13:bib:Bobojonov,zar14:bib:BobojonovA}]. 
 Также бурно развивается изучение сингулярных интегро-дифференциальных уравнений с~фредгольмовыми операторами в банаховых пространствах [\citen{zar15:bib:Falaleev,zar16:bib:FalaleevM,zar17:bib:FalaleevMB}]. 
 Но важно заметить, что в этих сингулярных интегро"=дифференциальных уравнениях существующие  интегралы понимаются в~смысле главного значения по Коши и поэтому для решения этих уравнений применяются методы теории аналитических функций. 
 В~случае, когда в~сингулярных интегро"=дифференциальных уравнениях интегралы понимаются в~обычном смысле Римана, т.~е. когда рассматривается уравнение вида
\begin{equation}
y'(x)+\frac{A}{x-a}y(x)+\int _{a}^{x}\frac{B}{(t-a)^{\alpha+1}}y(t)dt=f(x),
\label{zaripov:eq1}
\end{equation} 
где $\alpha=1$ или $\alpha>1$, исследователи не обращали особого внимания на решения таких уравнений. 
Это происходило потому, что ядро уравнений в~этих случаях будет нефредгольмовым и~поэтому изучение таких уравнений считалось неактуальным. 
Но оказывается, что в~зависимости от класса, в~котором ищутся решения уравнения \eqref{zaripov:eq1}, вопрос о~его изучении становится актуальным, причём даёт такие эффекты, которые в~существующей теории не наблюдались.

Заметим, что впервые в~монографии \cite{zar18:bib:Rajabov} была предложена новая методика изучения интегральных уравнений с~сингулярными и сверхсингулярными коэффициентами, которая потом в~работах [\citen{zar19:bib:RajabovN,zar20:bib:RajabovNR}] обобщалась на случай многомерных интегральных уравнений. 
Эта методика нами была использована для изучения интегро"=дифференциальных уравнений вида \eqref{zaripov:eq1}. 
Например, в~работах [\citen{zar21:bib:Zaripov,zar22:bib:ZaripovS}] исследовано уравнение вида \eqref{zaripov:eq1} в~случае, когда $\alpha=1$ или $\alpha>1$. 
Следующим этапом в~развитии теории уравнений вида  \eqref{zaripov:eq1} является случай, когда в~ядро этого уравнения добавляется логарифмическая особенность. 
Ниже мы будем исследовать именно такое  уравнение.

\smallskip

\Section[n]{Постановка задачи и ее решение}
Пусть $ \Gamma=\{x:a<x<b \}$ "--- множество точек на вещественной оси. 
На $ \Gamma$ рассмотрим модельное интегро"=дифференциальное уравнение вида \eqref{zaripov:eq1} с~логарифмической особенностью в~ядре при $\alpha=1$, т.~е. рассмотрим уравнение вида
\begin{equation}
y'(x)+\frac{P_{0}}{x-a}y(x)+\int _{a}^{x}
\Bigl[P_{1}+{P_{2}\,\ln\Bigl(\frac{x-a}{{t-a}}\Bigr)}\Bigr]
\frac{y(t)}{(t-a)^2}dt=f(x),
\label{zaripov:eq2}
\end{equation}
где $P_0$, $P_1$, $P_2 $ "--- заданные постоянные числа, $f(x)$ "--- заданная функция, $y(x)$ "--- искомая функция.

Прежде всего, через  $C'_a[a,b] $   обозначим класс таких функций, которые имеют производную первого порядка и~в~точке $x=a$  обращаются в~нуль с асимптотическим поведением
$$
y(x)=o[(x-a)^{\gamma_1}], \quad  \gamma_1>1.
$$

Решение однородного уравнения \eqref{zaripov:eq2} будем искать в классе $C'_a[a,b], $  т.~е. в виде $$y(x)=(x-a)^{\lambda}, \quad {\lambda}>1.$$ 
Тогда для определения $\lambda$  получим алгебраическое уравнение
\begin{equation}
\lambda+P_0+\frac{P_1}{\lambda-1}+\frac{P_2}{(\lambda-1)^2}=0,
\label{zaripov:eq3}
\end{equation}
которое назовем его характеристическим уравнением, соответствующим однородному уравнению \eqref{zaripov:eq2}. 
Исследуем уравнение  \eqref{zaripov:eq2} в зависимости от корней характеристического уравнения \eqref{zaripov:eq3}. 

\smallskip

\texttt{Случай I.1.}
Пусть корни характеристического уравнения \eqref{zaripov:eq3} являются вещественными и разными и они обозначены через $\lambda_1$, $\lambda_2$, $\lambda_3$. 
Кроме этого, пусть эти корни удовлетворяют условию $1<\lambda_1< \lambda_2< \lambda_3$, тогда, согласно общей теории, общее решение однородного уравнения \eqref{zaripov:eq2} даётся при помощи формулы 
$$
y(x)={(x-a)^{\lambda_1}}c_1+{(x-a)^{\lambda_2}}c_2+{(x-a)^{\lambda_3}}c_3,
$$
где $c_1$, $c_2$, $c_3$   "--- произвольные постоянные числа.

Используя метод вариации произвольных постоянных, решение неоднородного уравнения \eqref{zaripov:eq2} легко находим в таком виде:
 \begin{multline}
y(x)={(x-a)}^{{\lambda}_{1}}c_1+{(x-a)}^{{\lambda}_{1}}c_2+{(x-a)}^{{\lambda}_{1}}c_3-\\
-\frac{1}{{\Delta_0}P_2}\int _{a}^{x}
\Bigl[\Delta_1 (\lambda_1-1)^2\Bigl(\frac{x-a}{t-a}\Bigr)^{\lambda_1}+\\
+{\Delta_1} (\lambda_2-1)^2\Bigl(\frac{x-a}{t-a}\Bigr)^{\lambda_2}+
{\Delta_3}{({{\lambda_3}-1})^2}
\Bigl(\frac{x-a}{t-a}\Bigr)^{\lambda_3}\Bigr]f(t) dt,\label{zaripov:eq4}
\end{multline}
где
$$ \Delta_0=
\begin{vmatrix}
1 & 1 & 1 \\ {\lambda_1} & {\lambda_2} & {\lambda_3} \\ \frac{1}{{\lambda_1}-1} & \frac{1}{{\lambda_2}-1} & \frac{1}{{\lambda_3}-1}
\end{vmatrix},
$$
$\Delta_1$, $\Delta_2$, $\Delta_3$ "--- алгебраические дополнения элементов последней строки определителя $\Delta_0.$
 Для сходимости интегралов в правой части \eqref{zaripov:eq4} требуем, чтобы функция $f(x)$ в точке $x=a$ обращалась в нуль с асимптотическим поведением:
 \begin{equation}
f(x)=o\left[\left(x-a\right)^{\delta_1}\right], \quad  
{\delta_1}>{{\lambda_3}-1}.
\label{zaripov:eq5}
\end{equation}
Заметим, что если выполняется неравенство $1<\lambda_1< \lambda_2< \lambda_3$, то полученное решение \eqref{zaripov:eq4} принадлежит классу $C'_a[a,b] $.
Подставляя выражение \eqref{zaripov:eq4} в уравнение \eqref{zaripov:eq2}, легко находим, что оно в~действительности удовлетворяет этому уравнению.
Таким образом, доказана следующая теорема.

\smallskip

{\small \sc Теорема 1.} 
{\it Пусть в интегро-дифференциальном уравнении \eqref{zaripov:eq2} коэффициенты $P_0,$ $P_1,$  и $P_2$    такие$,$ что корни характеристического уравнения \eqref{zaripov:eq3} являются вещественными и~разными и~$1<\lambda_1< \lambda_2< \lambda_3.$ 
Функция $f(x)\in C[a,b]$   и  $f(a)=0$ с асимптотическим поведением \eqref{zaripov:eq5}$.$ 
Тогда однородное уравнение \eqref{zaripov:eq2} имеет три линейно независимых решения$,$ а неоднородное уравнение \eqref{zaripov:eq2} в~классе функций $y(x)\in C'_a[a,b],$ обращающихся в нуль в точке $x=a,$ всегда разрешимо и его общее решение содержит три произвольных постоянных и~даётся при помощи формулы} \eqref{zaripov:eq4}.

\smallskip

\texttt{Случай I.2.} 
Теперь пусть $\lambda_1<1< \lambda_2< \lambda_3$. 
В~этом случае из выражения \eqref{zaripov:eq4} следует, что если решение уравнения \eqref{zaripov:eq2} при $\lambda_1<1$  существует, то для того, чтобы оно принадлежало классу $ C'_a[a,b] $, произвольная постоянная $c_1$  должна быть равной нулю, т.~е. в~этом случае решение уравнения \eqref{zaripov:eq2} имеет вид
\begin{multline}
y(x)={(x-a)}^{{\lambda}_{1}}c_2+{(x-a)}^{{\lambda}_{1}}c_3
-\frac{1}{{\Delta_0}P_2}\int _{a}^{x}\Bigl[{\Delta_1}
{\left({{\lambda_1}-1}\right)^2}\Bigl(\frac{x-a}{t-a}\Bigr)^{\lambda_1}+\\
+{\Delta_1}
{\left({{\lambda_2}-1}\right)^2}\Bigl(\frac{x-a}{t-a}\Bigr)^{\lambda_2}+
{\Delta_3}{\left({{\lambda_3}-1}\right)^2}\Bigl(\frac{x-a}{t-a}\Bigr)^{\lambda_3}\Bigr]f(t)dt.\label{zaripov:eq6}
\end{multline}
В этом случае для сходимости интегралов в правой части \eqref{zaripov:eq6} функция $f(x)$ также должна обращаться в нуль с асимптотическим поведением \eqref{zaripov:eq5}.
Таким образом, имеет место следующая теорема.

\smallskip

{\small \sc Теорема 2.} {\it Пусть в интегро-дифференциальном уравнении \eqref{zaripov:eq2} коэффициенты $P_0,$ $P_1,$  и $P_2$  такие$,$ что корни характеристического уравнения \eqref{zaripov:eq3} являются вещественными и разными и $\lambda_1<1< \lambda_2< \lambda_3.$ 
Функция $f(x)\in C[a,b]$   и  $f(a)=0$ с асимптотическим поведением \eqref{zaripov:eq5}$.$ 
Тогда однородное уравнение \eqref{zaripov:eq2} имеет два линейно независимых решения$,$ а неоднородное уравнение \eqref{zaripov:eq2} в~классе функций $y(x)\in C'_a[a,b],$  обращающихся в нуль в точке $x=a,$  всегда разрешимо и его общее решение содержит две произвольные постоянные и даётся по формуле} \eqref{zaripov:eq6}.


\smallskip

\texttt{Случай I.3.}  Таким же образом доказывается, что  в~случае, когда $\lambda_1 \hm < \lambda_2<1< \lambda_3$ или $\lambda_1< \lambda_2< \lambda_3<1$б имеют место следующие теоремы.

\smallskip

{\small \sc Теорема 3.} {\it Пусть в интегро-дифференциальном уравнении \eqref{zaripov:eq2} коэффициенты $P_0,$ $P_1,$  и $P_2$  такие$,$ что корни характеристического уравнения \eqref{zaripov:eq3} являются вещественными и разными и $\lambda_1< \lambda_2<1< \lambda_3.$ 
Функция $f(x)\in C[a,b]$   и  $f(a)=0$ с асимптотическим поведением \eqref{zaripov:eq5}$.$ 
Тогда однородное уравнение \eqref{zaripov:eq2} имеет одно линейно независимое решение$,$ а неоднородное уравнение \eqref{zaripov:eq2} в классе функций $y(x)\in C'_a[a,b],$  обращающихся в нуль в точке $x=a,$  всегда разрешимо и~его общее решение содержит одну произвольную постоянную и~даётся по формуле}
\begin{multline}
y(x)={(x-a)}^{{\lambda}_{1}}c_3
-\frac{1}{{\Delta_0}P_2}\int _{a}^{x}
\Bigl[{\Delta_1}
{\left({{\lambda_1}-1}\right)^2}\Bigl(\frac{x-a}{t-a}\Bigr)^{\lambda_1}+\\
+{\Delta_1}
{\left({{\lambda_2}-1}\right)^2}\Bigl(\frac{x-a}{t-a}\Bigr)^{\lambda_2}+
{\Delta_3}{\left({{\lambda_3}-1}\right)^2}\Bigl(\frac{x-a}{t-a}\Bigr)^{\lambda_3}\Bigr]f(t)dt.\label{zaripov:eq7}
\end{multline}

\smallskip

{\small\sc Теорема 4.} {\it Пусть в интегро-дифференциальном уравнении \eqref{zaripov:eq2} коэффициенты $P_0,$ $P_1,$  и $P_2$  такие$,$ что корни характеристического уравнения \eqref{zaripov:eq3} являются вещественными и разными и $\lambda_1< \lambda_2< \lambda_3<1.$ Функция $f(x)\in C[a,b].$ 
Тогда однородное уравнение \eqref{zaripov:eq2} имеет только нулевое решение$,$ а неоднородное уравнение \eqref{zaripov:eq2} в классе функций $y(x)\in C'_a[a,b]$ имеет единственное решение$,$ которое даётся при помощи формулы}
\begin{multline}
y(x)=-\frac{1}{{\Delta_0}P_2}\int _{a}^{x}\Bigl[{\Delta_1}
{\left({{\lambda_1}-1}\right)^2}
\Bigl(\frac{x-a}{t-a}\Bigr)^{\lambda_1}+\\
+{\Delta_1}
{\left({{\lambda_2}-1}\right)^2}
\Bigl(\frac{x-a}{t-a}\Bigr)^{\lambda_2}+
{\Delta_3}{\left({{\lambda_3}-1}\right)^2}\Bigl(\frac{x-a}{t-a}\Bigr)^{\lambda_3}\Bigr]f(t)dt.\label{zaripov:eq8}
\end{multline}

\smallskip

Заметим, что в последнем случае для сходимости интеграла в правой части \eqref{zaripov:eq8} от функции $f(x)$ никаких условий кроме её непрерывности не требуется.

Таким образом, объединяя вышеприведенные результаты для интегро"=дифференциального уравнения \eqref{zaripov:eq2}, можно построить аналог теоремы об альтернативе Фредгольма в таком виде.

\smallskip

{\small \sc Аналог теоремы об альтернативе Фредгольма для интегро"=дифференциального уравнения \eqref{zaripov:eq2}, когда корни характеристического уравнения~\eqref{zaripov:eq3} являются вещественными и разными.} 
\emph{Если корни характеристического уравнения \eqref{zaripov:eq3} удовлетворяют условию $\lambda_1< \lambda_2< \lambda_3<1,$  то однородное уравнение \eqref{zaripov:eq2} имеет только тривиальное решение$,$ а неоднородное уравнение \eqref{zaripov:eq2} для каждой функции $f(x)\in C[a,b]$    имеет решение$,$ и притом единственное$,$ которое даётся по формуле \eqref{zaripov:eq8}$.$ 
Если корни характеристического уравнения \eqref{zaripov:eq3} удовлетворяют одному из условий $\lambda_1< \lambda_2<1< \lambda_3,$ $\lambda_1<1< \lambda_2< \lambda_3$ или $1<\lambda_1< \lambda_2< \lambda_3,$  то однородное уравнение \eqref{zaripov:eq2} имеет нетривиальные решения и его общее решение содержит либо одну$,$ либо две$,$ либо три  произвольных постоянных$,$ а неоднородное уравнение \eqref{zaripov:eq2} разрешимо тогда и только тогда$,$ когда его правая часть удовлетворяет условию \eqref{zaripov:eq5}$.$ 
В~этих случаях неоднородное уравнение \eqref{zaripov:eq2} имеет бесконечное число решений и~его общее решение содержит либо одну$,$ либо две$,$ либо три произвольных постоянных и~дается соответственно при помощи формулы \eqref{zaripov:eq7}$,$ \eqref{zaripov:eq6} или} \eqref{zaripov:eq4}.

\smallskip

\texttt{Случай II.1.} 
Пусть корни характеристического уравнения \eqref{zaripov:eq3} являются вещественными и равными $\lambda_1=\lambda_2=\lambda_3=\lambda$ и $\lambda>1$, тогда легко можно показать, что общее решение однородного уравнения \eqref{zaripov:eq2} задаётся формулой 
$$
y(x)={(x-a)^{\lambda}}c_4+{(x-a)^{\lambda}}\ln(x-a) c_5+{(x-a)^{\lambda}}\ln^2(x-a)c_6,
$$
где $c_4$, $c_5$, $c_6$   "--- произвольные постоянные числа.

 Используя метод вариации произвольных постоянных, после простых вычислений решение неоднородного уравнения \eqref{zaripov:eq2} находим в таком виде:
 \begin{multline}
y(x)={(x-a)^{\lambda}}c_4+{(x-a)^{\lambda}}\ln(x-a) c_5+{(x-a)^{\lambda}}\ln^2(x-a)c_6-\\
-\frac{(\lambda-1)^3}{2P_2}\int _{a}^{x}
\Bigl(\frac{x-a}{t-a}\Bigr)^\lambda
\Bigl[(\lambda-1)^2\ln^2\Bigl(\frac{x-a}{t-a}\Bigr)+\\
+4(\lambda-1)\ln
\Bigl(\frac{x-a}{t-a}\Bigr)+2\Bigr]f(t)dt.
\label{zaripov:eq9}
\end{multline}
Если поведение функции $f(x)$  в точке $x=a$  определяется из асимптотической формулы
\begin{equation}
f(x)=o\left[\left(x-a\right)^{\delta_2}\right], \quad  \delta_2>\lambda-1, \text{когда $x\rightarrow{a}$},\label{zaripov:eq10}
\end{equation}
то интеграл в правой части \eqref{zaripov:eq9} будет сходящимся.
Таким образом, имеет место следующая теорема.

\smallskip

{\small \sc Теорема 5.} {\it Пусть в интегро-дифференциальном уравнении \eqref{zaripov:eq9} коэффициенты $P_0,$ $P_1,$  и $P_2$  такие$,$ что корни характеристического уравнения \eqref{zaripov:eq3} являются вещественными и равными и $\lambda>1.$ 
Функция $f(x)\in C[a,b]$   и~${f(a)=0}$ с асимптотическим поведением \eqref{zaripov:eq10}$.$	
 Тогда однородное уравнение \eqref{zaripov:eq2} имеет три линейно независимых решения$,$ а неоднородное уравнение \eqref{zaripov:eq2} в классе функций $y(x)\in  C'_a[a,b]$ всегда разрешимо и его общее решение содержит три произвольных постоянных и даётся по формуле} \eqref{zaripov:eq9}.

\smallskip

\texttt{Случай II.2.} 
Если $\lambda<1$, то для того, чтобы решение вида \eqref{zaripov:eq9} принадлежало классу $ C'_a[a,b] $, произвольные постоянные $c_4$, $c_5$ и $c_6$  должны равняться нулю. 
Отсюда следует, что в~этом случае уравнение \eqref{zaripov:eq2} имеет только единственное решение и оно имеет вид
 \begin{multline}
y(x)=-\frac{(\lambda-1)^3}{2P_2}\int _{a}^{x} 
\Bigl(\frac{x-a}{t-a}\Bigr)^\lambda
\Bigl[(\lambda-1)^2\ln^2\Bigl(\frac{x-a}{t-a}\Bigr)+\\
+4(\lambda-1)\ln\Bigl(\frac{x-a}{t-a}\Bigr)+2\Bigr]f(t)dt.\label{zaripov:eq11}
\end{multline}
Таким образом, имеет место следующая теорема.

\smallskip

{\small \sc Теорема 6.} {\it Пусть в интегро-дифференциальном уравнении \eqref{zaripov:eq2} коэффициенты $P_0,$ $P_1,$  и $P_2$  такие$,$ что корни характеристического уравнения \eqref{zaripov:eq3} являются вещественными и равными и $\lambda<1.$ Функция $f(x)\in C[a,b].$
Тог\-да однородное уравнение \eqref{zaripov:eq2} имеет только нулевое решение$,$ а~неоднородное уравнение \eqref{zaripov:eq2} в~классе функций $y(x)\in C'_a[a,b]$ имеет единственное решение$,$ которое даётся формулой} \eqref{zaripov:eq11}.
 
 \smallskip

	Теперь построим аналог теоремы об альтернативе Фредгольма для уравнения \eqref{zaripov:eq2}.

 \smallskip

{\small\sc Аналог теоремы об альтернативе Фредгольма для интегро"=дифференциального уравнения \eqref{zaripov:eq2}, когда корни характеристического уравнения~\eqref{zaripov:eq3} являются вещественными и равными.} 
\emph{Если корни характеристического уравнения \eqref{zaripov:eq3} удовлетворяют условию $\lambda<1,$  то однородное уравнение \eqref{zaripov:eq2} имеет только тривиальное решение$,$ а~неоднородное уравнение \eqref{zaripov:eq2} для каждой функции $f(x)\in C[a,b]$  имеет единственное решение$,$ которое даётся формулой \eqref{zaripov:eq11}$.$ 
Если $\lambda>1,$  то однородное уравнение \eqref{zaripov:eq2} имеет три линейно независимых решения$,$ а~неоднородное уравнение \eqref{zaripov:eq2} разрешимо тогда и только тогда$,$ когда его правая часть удовлетворяет условию \eqref{zaripov:eq10}$.$ 
В~этом  случае неоднородное уравнение \eqref{zaripov:eq2} имеет бесконечное число решений$,$  его общее решение содержит три произвольных постоянных и даётся формулой}~\eqref{zaripov:eq9}.


\smallskip

\texttt{Случай III.1.}  
Пусть среди корней характеристического уравнения \eqref{zaripov:eq3} существуют комплексно"=сопряженные корни. 
Эти корни обозначим через ${\lambda_1=\lambda}$ и ${\lambda_{2,3}}\hm ={\alpha}\pm{i\beta}$. 
Если, не ограничивая общности, предположить, что $1<\lambda<\alpha$, то частными решениями однородного уравнения \eqref{zaripov:eq2} будут функции
 $$
 y_1(x)=(x-a)^{\lambda},\quad  y_2(x)=(x-a)^\alpha\cos\bigl[\beta\ln(x-a)\bigr],$$
 $$ y_3(x)=(x-a)^\alpha\sin\bigl[\beta\ln(x-a)\bigr]$$ 
 и его общее решение согласно общей теории можно записать в~таком виде:
  \begin{equation*}
y_1(x)=(x-a)^{\lambda}c_7+(x-a)^\alpha\cos\bigl[\beta\ln(x-a)\bigr]c_8+(x-a)^\alpha\sin\bigl[\beta\ln(x-a)\bigr]c_9.
\end{equation*}

Как и выше, используя метод вариации произвольных постоянных, решение неоднородного уравнения \eqref{zaripov:eq2} находим в виде
\begin{multline}
y(x)=(x-a)^{\lambda}c_7+(x-a)^\alpha\cos\bigl[\beta\ln(x-a)\bigr]c_8
+(x-a)^\alpha\sin\bigl[\beta\ln(x-a)\bigr]c_9-\\
-\frac{(\lambda-1)^2}{D_1P_2}\int _{a}^{x}
\Bigl(\frac{x-a}{t-a}\Bigr)^\lambda f(t)dt-
\frac{1}{D_1P_2\beta}\int _{a}^{x}
\Bigl(\frac{x-a}{t-a}\Bigr)^\alpha\times\\
\times\Bigl\{\alpha_1\sin\Bigl[\beta\ln\Bigl(\frac{x-a}{t-a}\Bigr)\Bigr]+
\beta_1\sin\Bigl[\beta\ln\Bigl(\frac{x-a}{t-a}\Bigr)\Bigr]\Bigr\}f(t)dt,
\label{zaripov:eq12}
\end{multline}
где $\alpha_1=(\alpha-1)^2(\alpha-\lambda)+\beta^2(\alpha+\beta-2)$, $\beta_1=(\alpha-1)\beta(2\lambda-\alpha-1)-\beta^3$.


Для сходимости интеграла в правой части \eqref{zaripov:eq12} от функции $f(x)$  потребуем, чтобы она в~точке $x=a$  обращалась в нуль с асимптотическим поведением
 {\begin{equation}
f(x)=o\left[\left(x-a\right)^{\delta_3}\right],\quad  \delta_3>\alpha-1, \text {когда $x\rightarrow{a}$}.\label{zaripov:eq13}
\end{equation}
Таким образом, имеет место следующая теорема.

\smallskip

{\small \sc Теорема 7.} {\it Пусть в интегро-дифференциальном уравнении \eqref{zaripov:eq2} коэффициенты $P_0,$ $P_1,$  и $P_2$  такие$,$ что среди корней характеристического уравнения \eqref{zaripov:eq3} существуют комплексно"=сопряженные корни и~выполняется условие $1<\lambda<\alpha.$ 
Функция $f(x)\in C[a,b]$   и~$f(a)=0$ с асимптотическим поведением~\eqref{zaripov:eq13}$.$
 Тогда однородное уравнение \eqref{zaripov:eq2} имеет три линейно независимых решения$,$ а~неоднородное уравнение \eqref{zaripov:eq2} в~классе функций $y(x)\in C'_a[a,b]$ всегда разрешимо и~его общее решение содержит три произвольных постоянных и~даётся при помощи формулы}~\eqref{zaripov:eq12}.

\smallskip

\texttt{Случай III.2.}  
Если выполняется условие $\lambda<1<\alpha$, то в общем решении \eqref{zaripov:eq12} произвольная постоянная $c_7$ должна быть равной нулю. В~этом случае решение уравнения \eqref{zaripov:eq2} имеет вид
\begin{multline}
y(x)=(x-a)^\alpha\cos
\bigl[\beta\ln(x-a)\bigr]c_8+(x-a)^\alpha\sin\bigl[\beta\ln(x-a)\bigr]c_9-\\
-\frac{(\lambda-1)^2}{D_1P_2}
\int _{a}^{x}\Bigl(\frac{x-a}{t-a}\Bigr)^\lambda f(t)dt-
\frac{1}{D_1P_2\beta}\int _{a}^{x}
\Bigl(\frac{x-a}{t-a}\Bigr)^\alpha\times\\
\times\Bigl\{\alpha_1\sin
\Bigl[\beta\ln\Bigl(\frac{x-a}{t-a}\Bigr)\Bigr]+
\beta_1\sin
\Bigl[\beta\ln\Bigl(\frac{x-a}{t-a}\Bigr)\Bigr]\Bigr\}f(t)dt.\label{zaripov:eq14}
\end{multline}

Очевидно, что при выполнении условия \eqref{zaripov:eq13} интегралы в правой части~\eqref{zaripov:eq14} будут сходящимися.
В этом случае имеет место следующая теорема.

\smallskip

{\small\sc Теорема 8.} {\it Пусть в интегро-дифференциальном уравнении \eqref{zaripov:eq2} коэффициенты $P_0,$ $ P_1,$  и $P_2$  такие$,$ что среди корней характеристического уравнения~\eqref{zaripov:eq3} существуют комплексно"=сопряженные корни и~выполняется условие $\lambda<1<\alpha.$ 
Функция $f(x)\in C[a,b]$   и~$f(a)=0$ с~асимптотическим поведением \eqref{zaripov:eq13}$.$
Тогда однородное уравнение \eqref{zaripov:eq2} имеет два линейно независимых решения$,$ а неоднородное уравнение \eqref{zaripov:eq2} в~классе функций $y(x)\in C'_a[a,b]$ всегда разрешимо и~его общее решение содержит две произвольных постоянных и~даётся при помощи формулы}~\eqref{zaripov:eq14}.

\smallskip

\texttt{Случай III.3.}  
В случае, когда выполняется условие $\lambda<\alpha<1$, в~общем решении \eqref{zaripov:eq12} все произвольные постоянные должны равняться нулю и~поэтому уравнение \eqref{zaripov:eq2} имеет единственное решение, которое даётся формулой
\begin{multline}
y(x)=-\frac{(\lambda-1)^2}{D_1P_2}\int _{a}^{x}
\Bigl(\frac{x-a}{t-a}\Bigr)^\lambda f(t)dt-
\frac{1}{D_1P_2\beta}\int _{a}^{x}
\Bigl(\frac{x-a}{t-a}\Bigr)^\alpha\times\\
\times
\Bigl\{\alpha_1\sin\Bigl[\beta\ln\Bigl(\frac{x-a}{t-a}\Bigr)\Bigr]+
\beta_1\sin
\Bigl[\beta\ln\Bigl(\frac{x-a}{t-a}\Bigr)\Bigr]\Bigr\}f(t)dt.
\label{zaripov:eq15}
\end{multline}

В этом случае от функции $f(x)$  никаких условий кроме её непрерывности не требуется.
Таким образом, имеет место следующая теорема.

\smallskip

{\small\sc Теорема 9.} {\it Пусть в интегро-дифференциальном уравнении \eqref{zaripov:eq2} коэффициенты $P_0,$ $ P_1,$  и $P_2$  такие$,$ что среди корней характеристического уравнения~\eqref{zaripov:eq3} существуют комплексно"=сопряженные корни и выполняется условие $\lambda<\alpha<1.$ 
Функция $f(x)\in C[a,b].$
 Тогда однородное уравнение \eqref{zaripov:eq2} имеет только тривиальное решение$,$ а неоднородное уравнение \eqref{zaripov:eq2} в классе функций $y(x)\in C'_a[a,b]$ имеет единственное решение$,$ которое даётся формулой}~\eqref{zaripov:eq15}.


\smallskip

Объединяя эти результаты, построим аналог теоремы об альтернативе Фредгольма для уравнения \eqref{zaripov:eq2}.

\smallskip

{\small \sc Аналог теоремы об альтернативе Фредгольма для интегро"=дифференциального уравнения \eqref{zaripov:eq2}, когда среди корней характеристического уравнения~\eqref{zaripov:eq3} существуют комплексно"=сопряженные корни.}
{\it  Если корни характеристического уравнения \eqref{zaripov:eq3} удовлетворяют условию $\lambda<\alpha<1,$ то однородное уравнение \eqref{zaripov:eq2} имеет только тривиальное решение$,$ а неоднородное уравнение \eqref{zaripov:eq2} для каждой функции $f(x)\in C[a,b]$  имеет единственное решение$,$ которое даётся формулой \eqref{zaripov:eq15}$.$ 
Если $\lambda<1<\alpha,$  то однородное уравнение \eqref{zaripov:eq2} имеет два линейно независимых решения$,$ а неоднородное уравнение \eqref{zaripov:eq2} разрешимо тогда и только тогда$,$ когда его правая часть удовлетворяет условию \eqref{zaripov:eq13}$.$ 
В этом  случае неоднородное уравнение \eqref{zaripov:eq2} имеет бесконечное число решений и~его общее решение содержит две произвольные постоянные и~выражается формулой~\eqref{zaripov:eq14}$.$ 
Если $1<\lambda<\alpha,$  то однородное уравнение \eqref{zaripov:eq2} имеет три линейно независимых решения$,$ а неоднородное уравнение \eqref{zaripov:eq2} разрешимо тогда и только тогда$,$ когда его правая часть удовлетворяет условию \eqref{zaripov:eq13}$.$ 
В этом  случае неоднородное уравнение \eqref{zaripov:eq2} имеет бесконечное число решений$,$  его общее решение содержит три произвольных постоянных и выражается формулой}~\eqref{zaripov:eq12}.

\smallskip

\Section[n]{Заключение}
В заключение заметим, что в общей теории число произвольных постоянных, входящих в решение интегро"=дифференциальных уравнений типа Вольтерра, равно наибольшему порядку производной, присутствующей в этих уравнениях, и эти  произвольные  постоянные легко  находятся с~помощью постановок различных начальных и~краевых задач. 
В~отличие от этого мы установили, что решение уравнения~\eqref{zaripov:eq2} в~общем случае содержит три произвольные постоянные. 
Также имеются случаи, когда уравнение \eqref{zaripov:eq2} содержит либо одну, либо две произвольные постоянные или выделяется случай, когда уравнение~\eqref{zaripov:eq2} имеет единственное решение без использования каких"=либо начальных или краевых задач. 
Отсюда вытекает, что только в~одном случае найденное нам решение согласуется с существующей теорией, а~присутствие других случаев показывает новизну данной работы. 
Далее эти эффекты дают нам возможность построить аналог теоремы Фредгольма для уравнения \eqref{zaripov:eq2}. Заметим, что раньше для интегро-дифференциального уравнения типа Вольтерра такие аналоги не  строились.

Другая особенность уравнения \eqref{zaripov:eq2} заключается в~том, что в~ядро этого уравнения добавляется логарифмическая особенность, за счёт которой число линейно независимых решений однородного уравнения \eqref{zaripov:eq2} не уменьшается, а~увеличивается.   	
Далее интегральные представления, полученные в данной статье, могут быть использованы для постановки и~решения различных начальных и краевых задач.

%	В следующих работах предполагается исследовать возможность построения аналога теоремы Фредгольма для интегро-дифференциального уравнения вида \eqref{zaripov:eq2} со сверхсингулярными коэффициентами, а в дальнейшем исследовать немодельные интегро-дифференциальные уравнения типа \eqref{zaripov:eq2}.

\AdditionalContent{Конкурирующие интересы}{Конкурирующих интересов не имею.} 
\AdditionalContent{Авторская ответственность}{Я несу полную ответственность за предоставление окончательной версии рукописи в печать. Окончательная версия рукописи мною одобрена.} 
\AdditionalContent{Место выполнения работы}{Работа выполнена на кафедре математического анализа и теории функций Таджикского национального университета.}
\AdditionalContent{Финансирование}{Исследование выполнялось без финансирования.}


\begin{thebibliography}{99}

\Bibitem{zar1:bib:Voltera}
\by Volterra~V.
\book Theory of functionals and of integral and integro-differential equations
\zmath{0086.10402}
\publaddr New York
\publ Dover Publ.
\totalpages 226 
\yr 1959

\RBibitem{zar2:bib:Magnaradze}
\paper Об одном новом интегральном уравнении теории крыла самолета
\by Магнарадзе~Л.~Г.
\jour Сообщ. АН ГрузССР
\yr 1942
\vol 3
\issue 6
\pages 503--508

\RBibitem{zar3:bib:Vekua}
\paper Об интегро-дифференциальном уравнении Прандтля
\by Векуа~И.~Н.
\jour Прикл. матем. мех.
\yr 1945
\vol 9
\issue 2
\pages 143--150

\RBibitem{zar4:bib:Veynberg}
\by Вайнберг~М.~М.
\paper Интегро-дифференциальные уравнения
\serial Итоги науки. Сер. Мат. анал. Теор. вероятн. Регулир. 1962
\yr 1964
\pages 5--37
\publ ВИНИТИ
\publaddr М.
\mathnet{http://mi.mathnet.ru/intv82}
\mathscinet{http://www.ams.org/mathscinet-getitem?mr=178320}
\zmath{https://zbmath.org/?q=an:0184.33702}

\RBibitem{zar5:bib:Nekrasov}
\by Некрасов~А.~И.
\paper Об одном классе линейных интегро"=дифференциальных уравнений
\jour  Тр. ЦАГИ
\yr  1934
\issue 190
\pages 1–25

\RBibitem{zar6:bib:Durdiev}
\paper Глобальная разрешимость одной обратной задачи для интегро-дифференциального уравнения электродинамики
\by Дурдиев~Д.~К.
\jour Дифференц. уравнения
\yr 2008
\vol 44
\issue 7
\pages 867--873
\elib{https://elibrary.ru/item.asp?id=11533343}

\RBibitem{zar7:bib:Durdiev}
\by Дурдиев~Д.~К.
\paper О единственности определения  ядра интегро-дифференциального уравнения параболического типа
\jour Вестн. Сам. гос. техн. ун-та. Сер. Физ.-мат. науки
\yr 2015
\vol 19
\issue 4
\pages 658--666
\mathnet{http://mi.mathnet.ru/vsgtu1444}
\crossref{10.14498/vsgtu1444}
\elib{http://elibrary.ru/item.asp?id=25687494}


\RBibitem{zar8:bib:Safarov}
\by Сафаров~Ж.~Ш.
\paper Оценки устойчивости решений некоторых обратных  задач для интегро-дифференциальных уравнений
\jour Вестн. Удмуртск. ун-та. Матем. Мех. Компьют. науки
\yr 2014
\issue 3
\pages 75--82
\mathnet{http://mi.mathnet.ru/vuu441}


\RBibitem{zar9:bib:Yuldoshev}
\by Юлдашев~Т.~К.
\paper Обратная задача для одного нелинейного интегро-дифференциального уравнения третьего порядка
\jour Вестн. СамГУ. Естественнонаучн. сер.
\yr 2013
\issue 9/1(110)
\pages 58--66
\mathnet{http://mi.mathnet.ru/vsgu400}

\RBibitem{zar10:bib:YuldoshevT}
\paper Об обратной задаче для нелинейных интегро-дифференциальных уравнений высшего порядка
\by Юлдашев~Т.~К.
\jour Вестник ВГУ. Серия: Физика. Математика
\yr 2014
\issue 1
\pages 153--163
\elib{https://elibrary.ru/item.asp?id=21445970}

\RBibitem{zar11:bib:YuldoshevT}
\by Юлдашев~Т.~К.
\paper Обратная задача для одного интегро-дифференциального уравнения Фредгольма  в~частных производных третьего порядка
\jour Вестн. Сам. гос. техн. ун-та. Сер. Физ.-мат. науки
\yr 2014
\issue 1(34)
\pages 56--65
\mathnet{http://mi.mathnet.ru/vsgtu1299}
\crossref{10.14498/vsgtu1299}
\elib{http://elibrary.ru/item.asp?id=22813960}


\RBibitem{zar12:bib:MagnaradzeG}
\paper Об одной системе линейных сингулярных интегро-дифференциальных уравнений и о линейной граничной задаче Римана
\by Магнарадзе~Л.~Г.
\jour Сообщ. АН ГрузССР
\yr 1943
\issue 5
\pages 3--9

\RBibitem{zar13:bib:Bobojonov}
\by Бободжанов~А.~А., Сафонов~В.~Ф.
\paper Метод нормальных форм в сингулярно возмущенных системах интегро-дифференциальных уравнений Фредгольма с быстро изменяющимися ядрами
\jour Матем. сб.
\yr 2013
\vol 204
\issue 7
\pages 47--70
\mathnet{http://mi.mathnet.ru/msb8139}
\crossref{10.4213/sm8139}
\mathscinet{http://www.ams.org/mathscinet-getitem?mr=3114874}
\zmath{https://zbmath.org/?q=an:06231583}
\adsnasa{http://adsabs.harvard.edu/cgi-bin/bib_query?2013SbMat.204..979B}
\elib{http://elibrary.ru/item.asp?id=20359261}
%\transl
%\by A.~A.~Bobodzhanov, V.~F.~Safonov
%\paper The method of normal forms for singularly perturbed systems of Fredholm integro-differential equations with rapidly varying kernels
%\jour Sb. Math.
%\yr 2013
%\vol 204
%\issue 7
%\pages 979--1002
%\crossref{10.1070/SM2013v204n07ABEH004327}
%\isi{http://gateway.isiknowledge.com/gateway/Gateway.cgi?GWVersion=2&SrcApp=PARTNER_APP&SrcAuth=LinksAMR&DestLinkType=FullRecord&DestApp=ALL_WOS&KeyUT=000324295300003}
%\elib{http://elibrary.ru/item.asp?id=21895806}
%\scopus{http://www.scopus.com/record/display.url?origin=inward&eid=2-s2.0-84888377177}

\RBibitem{zar14:bib:BobojonovA}
\by Бободжанов~А.~А., Сафонов~В.~Ф.
\paper ``Всплески'' в~интегро-дифференциальных уравнениях Фредгольма с~быстро изменяющимися ядрами
\jour Матем. заметки
\yr 2009
\vol 85
\issue 2
\pages 163--179
\mathnet{http://mi.mathnet.ru/mz4444}
\crossref{10.4213/mzm4444}
\mathscinet{http://www.ams.org/mathscinet-getitem?mr=2547997}
\zmath{https://zbmath.org/?q=an:1177.45009}
\elib{http://elibrary.ru/item.asp?id=13607171}
%\transl
%\jour Math. Notes
%\yr 2009
%\vol 85
%\issue 2
%\pages 153--167
%\crossref{10.1134/S0001434609010192}
%\isi{http://gateway.isiknowledge.com/gateway/Gateway.cgi?GWVersion=2&SrcApp=PARTNER_APP&SrcAuth=LinksAMR&DestLinkType=FullRecord&DestApp=ALL_WOS&KeyUT=000264327200019}
%\scopus{http://www.scopus.com/record/display.url?origin=inward&eid=2-s2.0-62949200440}

\RBibitem{zar15:bib:Falaleev}
\by Фалалеев~М.~В.
\paper Интегро-дифференциальные уравнения с фредгольмовым оператором при старшей производной в банаховых пространствах и их приложения
\jour Изв. Иркутского гос. ун-та. Сер. Математика
\yr 2012
\vol 5
\issue 2
\pages 90--102
\mathnet{http://mi.mathnet.ru/iigum70}

\RBibitem{zar16:bib:FalaleevM}
\by Фалалеев~М.~В.
\paper Сингулярные интегро-дифференциальные уравнения специального вида в банаховых пространствах и их приложения
\jour Изв. Иркутского гос. ун-та. Сер. Математика
\yr 2013
\vol 6
\issue 4
\pages 128--137
\mathnet{http://mi.mathnet.ru/iigum39}

\RBibitem{zar17:bib:FalaleevMB}
\by Фалалеев~М.~В.
\paper Вырожденные интегро-дифференциальные уравнения типа свертки в~банаховых пространствах
\jour Изв. Иркутского гос. ун-та. Сер. Математика
\yr 2016
\vol 17
\pages 77--85
\mathnet{http://mi.mathnet.ru/iigum274}

\RBibitem{zar18:bib:Rajabov}
\by Раджабов~Н.
\book Интегральные уравнения типов Вольтерра с фиксированными граничными и внутренними сингулярными и сверхсингулярными ядрами и их приложения
\publaddr Душанбе
\publ Деваштич
\yr 2007
\totalpages 221

\RBibitem{zar19:bib:RajabovN}
\by Раджабов~Н.
\paper Многомерное интегральное уравнение вольтеровского типа с сингулярными граничными областями в ядрах
\jour Докл. РАН
\yr 2011
\vol 437
\issue 2
\pages 158--161
\elib{https://elibrary.ru/item.asp?id=15639382}

\RBibitem{zar20:bib:RajabovNR}
\by Раджабов~Н., Раджабова~Л., Репин~О.~А.
\paper Об одном классе двумерных сопряженных интегральных уравнений вольтеровского типа
\jour Дифференц. уравнения
\yr 2011
\vol 47
\issue 9
\pages 1320--1330
\elib{https://elibrary.ru/item.asp?id=16766512}

\RBibitem{zar21:bib:Zaripov}
\paper Об одном классе модельного интегро-дифференциального уравнения первого порядка с одной сингулярной точкой в ядре
\by Зарипов~С.~К.
\jour Вестник Таджикского национального университета
\yr 2015
\issue 1/3(164)
\pages 27--32
\elib{https://elibrary.ru/item.asp?id=24732314}

\RBibitem{zar22:bib:ZaripovS}
\paper Об одном классе модельных интегро-дифференциальных уравнений первого порядка со сверх сингулярной точкой в ядре
\by Зарипов~С.~К.
\jour Вестник Таджикского национального университета
\yr 2015
\issue 1/6(191)
\pages 6--13
\elib{https://elibrary.ru/item.asp?id=25037274}




\end{thebibliography}


\newpage

\makeentitle

\AdditionalContent{Competing interests}{I have no competing interests.} 
\AdditionalContent{Author's Responsibilities}{I take full responsibility for submitting the final manuscript in print. I approved the final version of the manuscript.} 
\AdditionalContent{Plase of work}{The work was carried out at the Mathematical Analysis and Function Theory Department at the Tajik National University.}
\AdditionalContent{Funding}{The research has not had any sponsorship.}

\newpage


\begin{thebibliography}{99}

\Bibitem{zar1:bib:Voltera:en}
\by Volterra~V.
\book Theory of functionals and of integral and integro-differential equations
\zmath{0086.10402}
\publaddr New York
\publ Dover Publ.
\totalpages 226 
\yr 1959

\Bibitem{zar2:bib:Magnaradze:en}
\lang In Russian
\paper On a new integral equation of the theory of aircraft wings
\by Magnaradze~L.~G.
\jour Soobshch. AN GruzSSR
\yr 1942
\vol 3
\issue 6
\pages 503--508

\Bibitem{zar3:bib:Vekua:en}
\paper On the Prandtl integro-differential equation
\by Vekua~I.~N
\jour Prikl. matem. mekh.
\yr 1945
\vol 9
\issue 2
\pages 143--150
\elink \url{https://archive.org/details/nasa_techdoc_19880069126}

\Bibitem{zar4:bib:Veynberg:en}
\lang In Russian
\by Vainberg~M.~M.
\paper Integro-differential equations
\serial Itogi Nauki. Ser. Mat. Anal. Teor. Ver. Regulir. 1962
\yr 1964
\pages 5--37
\publ VINITI
\publaddr Moscow


\Bibitem{zar5:bib:Nekrasov:en}
\lang In Russian
\by Nekrasov~A.~I.
\paper On a Class of Linear Integro-Differential Equations
\jour  Tr. Tsentr. Aerogidrodin. Inst.
\yr  1934
\issue 190
\pages 1–25

\Bibitem{zar6:bib:Durdiev:en}
\paper Global solvability of an inverse problem for an integro-differential equation of electrodynamics
\by Durdiev~D.~K.
\jour Differ. Equ.
\yr  2008
\vol 44
\issue  7
\pages 893--899
\elib{https://elibrary.ru/item.asp?id=27875641}
\crossref{10.1134/S001226610807001X}

\Bibitem{zar7:bib:Durdiev:en}
\lang In Russian
\by Durdiev~D.~K.
\paper On the uniqueness of kernel determination in the integro-differential equation of parabolic type
\jour Vestn. Samar. Gos. Tekhn. Univ., Ser. Fiz.-Mat. Nauki \rm [J. Samara State Tech. Univ., Ser. Phys. Math. Sci.]
\yr 2015
\vol 19
\issue 4
\pages 658--666
\crossref{10.14498/vsgtu1444}



\Bibitem{zar8:bib:Safarov:en}
\lang In Russian
\by Safarov~Zh.~Sh.
\paper Evaluation of the stability of some inverse problems solutions for integro-differential equations
\jour Vestn. Udmurtsk. Univ. Mat. Mekh. Komp. Nauki
\yr 2014
\issue 3
\pages 75--82


\Bibitem{zar9:bib:Yuldoshev:w}
\lang In Russian
\by Yuldashev~T.~K.
\paper Inverse problem for a nonlinear integral and differential equation of the third order
\jour Vestnik Samarskogo Gosudarstvennogo Universiteta. Estestvenno-Nauchnaya Seriya
\yr 2013
\issue 9/1(110)
\pages 58--66

\Bibitem{zar10:bib:YuldoshevT:en}
\lang In Russian
\paper On inverse problem for nonlinear integro-differential equations of the higher order
\by Yuldashev~T.~K.
\jour Vestnik VGU. Seriia: Fizika. Matematika
\yr 2014
\issue 1
\pages 153--163


\Bibitem{zar11:bib:YuldoshevT:en}
\lang In Russian
\by Yuldashev~T.~K.
\paper Inverse Problem for a Fredholm Third Order Partial Integro-differential Equation
\jour Vestn. Samar. Gos. Tekhn. Univ., Ser. Fiz.-Mat. Nauki \rm [J. Samara State Tech. Univ., Ser. Phys. Math. Sci.]
\yr 2014
\vol 1(34)
\pages 56--65
\crossref{10.14498/vsgtu1299}



\Bibitem{zar12:bib:MagnaradzeG:en}
\lang In Russian
\paper On  a  system  of  linear singular integra-differential equations and on the 
linear Riemann boundary problem
\by Magnaradze~L.~G.
\jour Soobshch. AN GruzSSR
\yr 1943
\issue 5
\pages 3--9

\Bibitem{zar13:bib:Bobojonov:en}
\by Bobodzhanov~A.~A., Safonov~V.~F.
\paper The method of normal forms for singularly perturbed systems of Fredholm integro-differential equations with rapidly varying kernels
\jour Sb. Math.
\yr 2013
\vol 204
\issue 7
\pages 979--1002
\crossref{10.1070/SM2013v204n07ABEH004327}
\isi{http://gateway.isiknowledge.com/gateway/Gateway.cgi?GWVersion=2&SrcApp=PARTNER_APP&SrcAuth=LinksAMR&DestLinkType=FullRecord&DestApp=ALL_WOS&KeyUT=000324295300003}
\elib{http://elibrary.ru/item.asp?id=21895806}
\scopus{http://www.scopus.com/record/display.url?origin=inward&eid=2-s2.0-84888377177}

\Bibitem{zar14:bib:BobojonovA:en}
\paper ``Splashes'' in Fredholm integro-differential equations with rapidly varying kernels
\by  Bobodzhanov~A. A.,  Safonov~V.~F.
\jour Math. Notes
\yr 2009
\vol 85
\issue 2
\pages 153--167
\crossref{10.1134/S0001434609010192}
\isi{http://gateway.isiknowledge.com/gateway/Gateway.cgi?GWVersion=2&SrcApp=PARTNER_APP&SrcAuth=LinksAMR&DestLinkType=FullRecord&DestApp=ALL_WOS&KeyUT=000264327200019}
\scopus{http://www.scopus.com/record/display.url?origin=inward&eid=2-s2.0-62949200440}

\Bibitem{zar15:bib:Falaleev:en}
\lang In Russian
\by Falaleev~M.~V.
\paper Integro-differential equations with Fredholm operator by the derivative of the higest order in Banach spaces and it's applications
\jour IIGU Ser. Matematika
\yr 2012
\vol 5
\issue 2
\pages 90--102

\Bibitem{zar16:bib:FalaleevM:en}
\lang In Russian
\by Falaleev~M.~V.
\paper Singular integro-differential equations of the special type in Banach spaces and it's applications
\jour IIGU Ser. Matematika
\yr 2013
\vol 6
\issue 4
\pages 128--137

\Bibitem{zar17:bib:FalaleevMB:en}
\lang In Russian
\by Falaleev~M.~V.
\paper Degenerate integro-differential equations of convolution type in Banach spaces
\jour IIGU Ser. Matematika
\yr 2016
\vol 17
\pages 77--85

\Bibitem{zar18:bib:Rajabov:en}
\lang In Russian
\by Rajabov~N.
\book Integral'nye uravneniia tipov Vol'terra s fiksirovannymi granichnymi i vnutrennimi singuliarnymi i sverkhsinguliarnymi iadrami i ikh prilozheniia \rm [Integral equations of Volterra types with fixed boundary and internal singular and supersingular kernels and their applications]
\publaddr Dushanbe
\publ Devashtich
\yr 2007
\totalpages 221


\Bibitem{zar19:bib:RajabovN:en}
\paper Multidimensional Volterra-type integral equation with singular boundary domains in kernels
\by Rajabov~N.
\jour Dokl. Math. 
\yr  2011
\vol  83
\issue 2
\pages 165--168
\elib{https://elibrary.ru/item.asp?id=27914371}
\crossref{10.1134/S1064562411020116}

\Bibitem{zar20:bib:RajabovNR:en}
\paper On a class of two-dimensional adjoint integral equations of Volterra type
\by Rajabova~L., Rajabov~N., Repin~O.~A.
\jour Differ. Equ.
\yr 2011
\vol 47
\issue 9
\pages 1333--1343
\elib{https://elibrary.ru/item.asp?id=18014291}
\crossref{10.1134/S0012266111090102}

\Bibitem{zar21:bib:Zaripov:en}
\lang In Russian
\paper On a class of model first-order integro-differential equations with one singular point in the kernel
\by Zaripov~S.~K.
\jour Vestnik Tadzhikskogo natsional'nogo universiteta
\yr 2015
\issue 1/3(164)
\pages 27--32


\Bibitem{zar22:bib:ZaripovS:en}
\lang In Russian
\paper On a class of model first-order integro-differential equations with an additional  singular point in the kernel
\by Zaripov~S.~K.
\jour Vestnik Tadzhikskogo natsional'nogo universiteta
\yr 2015
\issue 1/6(191)
\pages 6--13




\end{thebibliography}


\renewcommand{\baselinestretch}{0.9}

%%
%\end{document}


